\documentclass[en,hazy,blue,12pt,normal]{elegantnote}

\input{preface}

% (If your preface already loads these, it's still usually fine.)
\usepackage{amsmath, amssymb}
\usepackage{graphicx}
\usepackage{booktabs}

\title{Winter 2026 DSC 240: Homework 3}
\author{Zeyu Bian}
\date{\today}

\begin{document}

\maketitle
\begin{homework}
    (10 points): Maximum Likelihood Estimation for Bernoulli Distribution

Let $X_1, X_2, \ldots, X_n$ be $n$ independent and identically distributed (i.i.d.) random variables following a Bernoulli distribution. The Bernoulli distribution for a random variable $X$ is given by:

$$
\mathbb{P}(X=x ; p)=p^x(1-p)^{1-x}
$$

where $x \in\{0,1\}$ and $p$ is the probability of $X$ being 1 . Derive the maximum likelihood estimator (MLE) for the parameter $p$ using the given samples.
\end{homework}

\begin{proof}
    We directly calculate the log-likelihood:

$$
\ell(p)=\log L(p)=\sum_{i=1}^n\left[x_i \log p+\left(1-x_i\right) \log (1-p)\right] .
$$

Take the derivative w.r.t to p and let it to be zero:
\[\frac{d \ell}{d p}=\sum_{i=1}^n\left(\frac{x_i}{p}-\frac{1-x_i}{1-p}\right)=0.\]

We get the MLE estimator
\[\hat{p} = \frac{\sum_{i=1}^nx_i}{n}.\]
\end{proof}

\begin{homework}
    ( 20 points): Cross-Entropy Loss and Logistic Loss
Consider a binary classification problem with outputs $y \in\{0,1\}$. Assume $p(y \mid x)$ is $\operatorname{Bernoulli}\left(\hat{p}\left(x ; \theta^*\right)\right)$ for some parameter vector $\theta^* \in \mathbb{R}^d$, where for each $\theta, \hat{p}(x ; \theta)$ returns a number between $[0,1]$.

(a) ( $\mathbf{1 0}$ points) For each parameter $\theta$, the cross-entropy loss for a set of $n$ examples with predictions $\hat{p}_i:=\hat{p}\left(x_i ; \theta\right)$ and true labels $y_i$ is given by:

$$
\operatorname{CE}(\theta)=-\frac{1}{n} \sum_{i=1}^n\left(y_i \log \left(\hat{p}_i\right)+\left(1-y_i\right) \log \left(1-\hat{p}_i\right)\right) .
$$


Derive this loss function starting from the principle of Maximum Likelihood Estimation (MLE).

(b) ( $\mathbf{1 0}$ points) Instantiate the classifier's prediction $\hat{p}(x ; \theta)$ as the sigmoid function applied to a linear combination of input features, i.e., $\hat{p}(x ; \theta)=\sigma\left(\theta^{\top} x\right)$, where $\sigma(z)=\frac{1}{1+e^{-z}}$. Show that in this case, minimizing the cross-entropy loss above is equivalent to minimizing the familiar logistic loss, which is given by:

$$
L(\theta)=\frac{1}{n} \sum_{i=1}^n \log \left(1+e^{-\tilde{y}_i \theta^{\top} x_i}\right) \quad \text { where } \tilde{y}_i=\left\{\begin{array}{ll}
1 & \text { if } y_i=1 \\
-1 & \text { if } y_i=0
\end{array} .\right.
$$
\end{homework}

\begin{proof}
(a)
The likelihood over the dataset is
    \[L(\theta)=\prod_{i=1}^n \hat{p}_i^{y_i}\left(1-\hat{p}_i\right)^{1-y_i}.\]

    Maximize log-likelihood $\Leftrightarrow$ minimize negative average log-likelihood:

$$
-\frac{1}{n} \log L(\theta)=-\frac{1}{n} \sum_{i=1}^n\left(y_i \log \hat{p}_i+\left(1-y_i\right) \log \left(1-\hat{p}_i\right)\right),
$$

which is exactly the cross-entropy expression.

(b)Plug in the expression of $\hat{p}$ and unify the case for $y_i$, we have:
\begin{align}
    \textbf{CE}(\theta)&=\frac{1}{n}\sum_{i=1}^n(\log(1+e^{-\theta^\top x})I_{y_i = 1} + \log(1+e^{\theta^\top x})I_{y_i = 0})\\
    &=\frac{1}{n}\sum_{i=1}^n  \log \left(1+e^{-\bar{y}_i \theta^{\top} x_i}\right)\\
    &= L(\theta).
\end{align}
\end{proof}
\begin{homework}
    $(40+10$ bonus points): Shape of Bayes Optimal Classifiers for Gaussian Data

Assume we have a binary classification problem where:

- Data points from class A are generated from a Gaussian distribution $\mathcal{N}\left(\mu_1, I_d\right)$.

- Data points from class B are generated from a Gaussian distribution $\mathcal{N}\left(\mu_2, I_d\right)$.

- The prior probabilities for class A and class B are equal, $P$ (class A$)=P($ class B$)=0.5$.

(a) ( $\mathbf{1 0}$ points) Derive the posterior probability $P(y \mid x)$ for class A and B and the decision rule for the Bayes optimal classifier. Given $\mu_1=[0,0]$ and $\mu_2=[1,1]$, plot the decision boundary and determine if it is linear.
(Hint: Recall that the Bayes optimal classifier $h^*$ is given by $h^*(x)=\arg \max _y p(y \mid x)$. So the strategies for solving Part (a) and (b)(c)(d) is to first derive $p(y \mid x)$ using the Bayes rule. )

(b) ( $\mathbf{1 0}$ points) Consider the case where the prior probabilities are not equal, $P($ class A$)=p$ and $P($ class B$)=1-p$. Derive the Bayes optimal classifier. For the 2D example, let $\mu_1=[0,0]^T$, $\mu_2=[1,1]^T$, and $p=0.8$. Plot the decision boundary.

(c) ( $\mathbf{1 0}$ points) Now let the covariance matrix for class A be $\sigma^2 I_d$ while keeping class B's covariance matrix as $I_d$. Derive the Bayes optimal classifier. For the 2D example, let $\mu_1= [0,0]^T, \mu_2=[1,1]^T, p=0.5$, and $\sigma^2=0.25$. Plot the decision boundary.

(d) ( $\mathbf{1 0}$ points) Instantiate the decision rule from (c) with $\mu_1=\mu_2$ what is the shape of the decision boundary? For the scenario where $\mu_1=[0,0]^T$ and $\mu_2=[0,0]^T$ with $p=0.5$ and $\sigma^2=0.25$, plot the decision boundary.

(e) (Bonus 10 points) Examine whether the decision boundaries from parts (a)-(d) can be realized by a Gaussian Naive Bayes model. Provide a reasoning for your answer. (Hint: Whether the Bayes Optimal classifier can be realized by $\mathcal{H}$ means whether $h^* \in \mathcal{H}$, i.e., whether there is a parameter vector of the Gaussian Naive Bayes classifier that is the same as $h^*$ )

Deliverables for each subproblem (a) (b) (c) and (d):

1. Analytical expressions for $P(y \mid x)$.

2. The decision rule of the Bayes optimal classifier / analytical expression of the decision boundary.

3. Plotted figures that clearly demonstrate the decision boundary along with 100 random samples (color-coded in red and blue) from each class in the same figure generated in the specified way.

\end{homework}

\begin{proof}[\solutionname]
    Bayes rule:
    \[P(A \mid x)=\frac{\pi_1 p(x \mid A)}{\pi_1 p(x \mid A)+\pi_2 p(x \mid B)}, \quad P(B \mid x)=1-P(A \mid x) .\]

and the optimal solution 
\[h^*(x)=\arg \max _{y \in\{A, B\}} P(y \mid x) .\]

    (a)For Gaussian, we have 
    \[\log p(x \mid y)=-\frac{d}{2} \log (2 \pi)-\frac{1}{2} \log \left|\Sigma_y\right|-\frac{1}{2}\left(x-\mu_y\right)^{\top} \Sigma_y^{-1}\left(x-\mu_y\right) .\]
    Therefore, the equivalent discriminant rule:
    
Choose A if $\log \pi_1+\log p(x \mid A)>\log \pi_2+\log p(x \mid B)$.

    Since $\pi_1 = \pi_2, \Sigma_A = \Sigma_B$, we see the quadratic terms and prior terms cancel out. Therefore the decision boundry is linear.

    \[p(A|x) = \frac{1}{1+e^{x^\top(\mu_B-\mu_A)+\frac{1}{2}(\|\mu_A\|^2 - \|\mu_B\|^2)}}, p(B|x) \text{ is obvious by symmetry.}\]

    \begin{figure}[!htbp]
        \centering
        \includegraphics[width=\linewidth]{Figures/hw3/q3a_boundary.png}
    \end{figure}
    
(b) 
Decide A if $\delta(x)>0$, else decide B, where $\delta(x)=\left(\mu_1-\mu_2\right)^{\top} x-\frac{1}{2}\left(\left\|\mu_1\right\|^2-\left\|\mu_2\right\|^2\right)+\log \frac{p}{1-p}$.

\begin{figure}[!htbp]
    \centering
    \includegraphics[width=\linewidth]{Figures/hw3/q3b_boundary.png}
\end{figure}

(c)optimal classifier:
\[\left\|x-\mu_2\right\|^2-\frac{1}{\sigma^2}\left\|x-\mu_1\right\|^2+2 \log \frac{p}{1-p}-d \log \left(\sigma^2\right)>0 \Rightarrow \operatorname{decide} A,\]
\begin{figure}[!htbp]
    \centering
    \includegraphics[width=\linewidth]{Figures/hw3/q3c_boundary.png}
\end{figure}

(d) The shape is an ellipse. 
\begin{figure}[!htbp]
    \centering
    \includegraphics[width=\linewidth]{Figures/hw3/q3d_boundary.png}
\end{figure}

(e)
They all can. Because GNB assumes conditional independence of features given the class, yielding a quadratic boundary in general, and a hypersphere when $\mu_1 =\mu_2$. For (a), (b), to force a linear boundary in a “naive” diagonal-covariance style, we can  tie/pooled variances across classes or use a model that assumes shared covariance (e.g., LDA / logistic regression).

\end{proof}
\newpage

\begin{homework}
    (30 points): Naive Bayes Classifier vs Linear Classifier
In this question, we have provided you with three carefully designed binary classification datasets (data-1.npy, data-2.npy, and data-3.npy).

(a) ( $\mathbf{1 5}$ points) For each of the three datasets provided, fit a Gaussian Naive Bayes classifier and a linear logistic regression classifier using the sklearn library in Python. Plot the dataset with points color-coded according to their true labels. Additionally, plot the decision boundaries for both classifiers.
(Hint: You may use matplotlib for plotting and sklearn's GaussianNB and LogisticRegression for model fitting. The first two columns in the dataset represents the features and the last column is labels.)

(b) ( $\mathbf{1 5}$ points) Discuss why the Naive Bayes classifier and / or the linear classifier fail on each example. Justify your answers.
\end{homework}

\begin{proof}[\solutionname]
(a) \begin{figure}[!htbp]
    \centering
    \includegraphics[width=\linewidth]{Figures/hw3/q4_data1_boundaries.png}
\end{figure}

\begin{figure}[!htbp]
    \centering
    \includegraphics[width=\linewidth]{Figures/hw3/q4_data2_boundaries.png}
\end{figure}

\begin{figure}[!htbp]
    \centering
    \includegraphics[width=\linewidth]{Figures/hw3/q4_data3_boundaries.png}
\end{figure}

    (b)data-1 (two lines; LR better): The class boundary is essentially linear. Logistic Regression fits a single straight separating hyperplane, which matches the geometry. GaussianNB assumes each class is Gaussian; if the data lie near thin lines, the Gaussian fit is a poor model, so its boundary can be biased.
    
data-2 (two ellipses; GNB better): The classes look like (possibly rotated) Gaussian clusters. GaussianNB models each class with a Gaussian and yields a curved/quadratic boundary, which fits elliptical structure (but here the ellipse is skew, meaning correlation between features, so GNB is not that good somewhow). Logistic Regression is linear, so it cannot separate two ellipses well and underfits.

data-3 (two circles; GNB better): The class boundary is nonlinear (roughly circular). Logistic Regression's linear boundary fails. GaussianNB still works because each class is approximately Gaussian around its center, producing a quadratic boundary that can match circular separation.
\end{proof}
\end{document}